% Pregunta 886 (2003-julio-OEP2003, n.º 14): el progreso técnico desplaza hacia arriba la función de producción; con el mismo k se obtiene más producto y la pendiente (productividad marginal) es mayor. Autor: Víctor Gutiérrez Marcos
\documentclass[tikz,border=6pt]{standalone}
\input{tikz-tcee}
\begin{document}
\begin{tikzpicture}[tcee,x=1cm,y=0.9cm]
  % f0(k) = 1.5 sqrt(k), f1(k) = 2.2 sqrt(k); k0 = 2.5
  % f0(k0) = 2.372, f0'(k0) = 0.474 ; f1(k0) = 3.479, f1'(k0) = 0.696
  \ejes{$k$}{$y$}{6.6}{5.8}
  \draw[curva1,domain=0:6.2,samples=80] plot (\x,{1.5*sqrt(\x)});
  \node[etiqueta,tceemorado,anchor=west] at (6.2,3.73) {$f_0(k)$};
  \draw[curvanueva,desplazada,domain=0:6.2,samples=80] plot (\x,{2.2*sqrt(\x)});
  \node[etiqueta,tceerojo,anchor=west] at (6.2,5.48) {$f_1(k)$};
  % Tangentes en k0
  \draw[tceemorado,thick] (1.1,1.708) -- (4.3,3.226);
  \draw[tceerojo,thick] (1.1,2.505) -- (4.3,4.732);
  \draw[guia] (2.5,0) -- (2.5,3.479);
  \draw[guia] (0,2.372) -- (2.5,2.372);
  \draw[guia] (0,3.479) -- (2.5,3.479);
  \node[punto] at (2.5,2.372) {};
  \node[punto] at (2.5,3.479) {};
  \node[below] at (2.5,0) {$k_0$};
  \node[left] at (0,2.372) {$y_0$};
  \node[left] at (0,3.479) {$y_1$};
  \draw[flecha,tceerojo] (2.75,2.47) -- (2.75,3.38);
  \node[etiqueta,font=\footnotesize,align=left,anchor=west] at (3.3,1.2) {\textcolor{tceerojo}{tangente roja} más inclinada\\que la \textcolor{tceemorado}{morada}: mayor\\productividad marginal};
\end{tikzpicture}
\end{document}
